
\section{OLD Particle Geometry} \label{Fundamental}

We need to have a relationship between velocity in ($v_{in}$) (just prior to the collision), the velocity at any degree of penetration ($v_p$), the total circle area ($A_T$) and penetration area ($A_P$). In an \textit{\textbf{instantaneous, purely geometric}} system is that the ratio of incoming velocity (the velocity just before the collision occurs) to penetration velocity is proportional to the ratio of total area to penetration area.
\begin{equation}
\begin{aligned}
\frac{v_{in}}{v_p} \propto \frac{A_T}{A_P}
\end{aligned}
\end{equation}


The proportionality factor is the force with which the particle resists collision at depth of penetration ($\alpha$) or repulsion force ($F_R(\alpha)$) being an increasing parabolic function. So,

\begin{equation}
\begin{aligned}
\frac{v_{in}}{v_p} = F_R(\alpha) \frac{A_T}{A_P}
\end{aligned}
\end{equation}

and,

\begin{equation}
\begin{aligned}
v_{in}\frac{1}{F_R(\alpha)}\frac{A_P}{A_T} =  {v_p}
\end{aligned}
\end{equation}


\section{A second collision function}
A second collision function is that the change in velocity is proportional to the penetration area,

\begin{equation}
\begin{aligned}
(v_{in} - v_{p}) = P A_P
\end{aligned}
\end{equation}


So that 

\begin{equation}
\begin{aligned}
v_{p} = v_{in}-P A_P
\end{aligned}
\end{equation}

At the turnaround point, the penetration velocity is zero:

\begin{equation}
\begin{aligned}
0 = v_{in}-P A_0
\end{aligned}
\end{equation}

so the zero-velocity penetration area can be predicted from incoming velocity:

\begin{equation}
\begin{aligned}
A_0 = \frac{v_{in}}{P}
\end{aligned}
\end{equation}

\section{Equivalence of the two functions}

This section derives the algebraic force factor needed to make the first
collision function and second collision function describe the same velocity
curve.

The first function gives:

\begin{equation}
\begin{aligned}
v_p = v_{in}\frac{1}{F_R(\alpha)}\frac{A_P}{A_T}
\end{aligned}
\end{equation}

The second function gives:

\begin{equation}
\begin{aligned}
v_p = v_{in}-P A_P
\end{aligned}
\end{equation}

For the two functions to be equivalent, the two expressions for $v_p$ must be equal:

\begin{equation}
\begin{aligned}
v_{in}\frac{1}{F_R(\alpha)}\frac{A_P}{A_T}
=
v_{in}-P A_P
\end{aligned}
\end{equation}

Solving for the resistance function gives:

\begin{equation}
\begin{aligned}
F_R(\alpha)
=
\frac{v_{in}A_P}
{A_T(v_{in}-P A_P)}
\end{aligned}
\end{equation}

Using the turnaround relation $A_0 = v_{in}/P$, this can also be written as:

\begin{equation}
\begin{aligned}
F_R(\alpha)
=
\frac{A_P A_0}
{A_T(A_0-A_P)}
\end{aligned}
\end{equation}

Thus the first function and second function can describe the same collision
curve if $F_R(\alpha)$ is chosen consistently.  As $A_P$ approaches $A_0$, the
denominator approaches zero, so $F_R(\alpha)$ grows without bound.  This
corresponds to the zero-velocity turnaround point.


\section{Force factor}


This section gives a physical force law that can be used to compute the
turnaround depth by energy balance.  This is different from the algebraic
equivalence factor in Section 3.  The Section 3 factor is the exact function
needed to make two velocity equations equivalent.  The force law here describes
how collision resistance grows with penetration depth.

Assume the repulsive force is an increasing parabolic function of penetration
depth:
\begin{equation}
\begin{aligned}
F(\alpha) = Q\alpha^2
\end{aligned}
\end{equation}

where $Q$ is the stiffness factor.

At the zero-velocity turnaround point, incoming normal kinetic energy has been
stored by the repulsive force:

\begin{equation}
\begin{aligned}
\frac{1}{2}\mu v_{in}^2
=
\int_0^{\alpha_0} Q\alpha^2\,d\alpha
\end{aligned}
\end{equation}

Evaluating the integral:

\begin{equation}
\begin{aligned}
\frac{1}{2}\mu v_{in}^2
=
\frac{Q\alpha_0^3}{3}
\end{aligned}
\end{equation}

Solving for turnaround penetration depth:

\begin{equation}
\begin{aligned}
\alpha_0
=
\left(\frac{3\mu v_{in}^2}{2Q}\right)^{1/3}
\end{aligned}
\end{equation}

The zero-velocity penetration area is then:

\begin{equation}
\begin{aligned}
A_0 = A(\alpha_0)
\end{aligned}
\end{equation}

Thus Section 3 gives the algebraic force factor for equivalence, while this
section gives a physical force law for predicting the turnaround depth.

\section{Impact on starting in a collision}

When particles start in a collision, the configured velocity is interpreted as
the incoming velocity before turnaround.  The current overlap is assumed to be
on the rebound side of the collision history.

With the physical force law, the zero-velocity penetration depth is no longer
only a configured fraction.  It can be predicted from the incoming velocity and
the stiffness factor $Q$:

\begin{equation}
\begin{aligned}
\alpha_0
=
\left(\frac{3\mu v_{in}^2}{2Q}\right)^{1/3}
\end{aligned}
\end{equation}

and:

\begin{equation}
\begin{aligned}
A_0 = A(\alpha_0)
\end{aligned}
\end{equation}

For a starting overlap with current penetration area $A_P$:

\begin{itemize}[leftmargin=*]
    \item If $A_P < A_0$, the starting state is physically valid and is on the
    rebound path.
    \item If $A_P = A_0$, the starting state is exactly at the zero-velocity
    turnaround point.
    \item If $A_P > A_0$, the starting overlap is deeper than the incoming
    velocity and stiffness factor could physically produce.
\end{itemize}

The last case requires an explicit policy.  The simulation can:

\begin{itemize}[leftmargin=*]
    \item promote $A_0$ to the current penetration area, as the current Python
    model does,
    \item reject the configuration as physically impossible,
    \item or infer the larger incoming velocity required to reach the current
    penetration area.
\end{itemize}

Thus, for starting collisions, the force-law model replaces an arbitrary
\texttt{zero\_velocity\_overlap\_fraction} with a predicted turnaround point
based on incoming velocity, reduced mass, and stiffness.